\begin{definitionbox}{Definition (Limit Point)}
\begin{definition}[Limit Point]
\label{def:limit-point}
Let $A\subseteq\mathbb{R}$ and $x\in\mathbb{R}$.  The point $x$ is a
\emph{limit point} of $A$ if every deleted open ball centered at $x$ contains
a point of $A$; that is,
\[
  (\forall r>0)\bigl((B_r(x)\setminus\{x\})\cap A\neq\varnothing\bigr).
\]
Equivalently,
\[
  (\forall r>0)(\exists a\in A)(0<|a-x|<r).
\]
\end{definition}
\end{definitionbox}
\begin{remark*}[Standard quantified statement]
\[
x \text{ is a limit point of } E \Longleftrightarrow \forall r>0\;\exists y\in E\;(0<|y-x|<r).
\]
\end{remark*}

\begin{remark*}[Predicate reading]
\[
E\subseteq\mathbb{R},\qquad x\in\mathbb{R},\qquad
\operatorname{IsCluster}(x,E,\mathbb{R})
\Longleftrightarrow
\forall r>0\;\exists y\in E\;(0<|y-x|<r).
\]
\end{remark*}
\begin{remark*}[Negated quantified statement]
\[
x \text{ is not a limit point of } E \Longleftrightarrow \exists r>0\;\forall y\in E\;(y=x \lor |y-x|\geq r).
\]
\end{remark*}

\begin{remark*}[Negation predicate reading]
\[
x \text{ is not a limit point of } E \Longleftrightarrow \exists r>0\;\forall y\in E\;(y=x \lor |y-x|\geq r).
\]

The negation predicate reading is the predicate-level form of the displayed negated statement.
\end{remark*}

\begin{remark*}[Failure modes]
\begin{description}
  \item[Exposition.]
  Failure is witnessed by the displayed negated or contrapositive condition.

  \item[\operatorname{IsCluster}.]
  \[
x \text{ is not a limit point of } E \Longleftrightarrow \exists r>0\;\forall y\in E\;(y=x \lor |y-x|\geq r).
\]
\[
E\subseteq\mathbb{R},\qquad x\in\mathbb{R},\qquad
\operatorname{IsCluster}(x,E,\mathbb{R})
\Longleftrightarrow
\forall r>0\;\exists y\in E\;(0<|y-x|<r).
\]
\end{description}
\end{remark*}
\begin{remark*}[Interpretation]
Every punctured ball around a limit point meets $E$: the set accumulates at $x$. The strict inequality $0<|y-x|$ excludes $x$ itself, so $x$ need not belong to $E$ to be a limit point. This is the exact condition under which a sequence in $E$ can converge to $x$.
\end{remark*}
\begin{dependencies}
\begin{itemize}
  \item \hyperref[def:open-ball]{Open Ball on the Real Line}
\end{itemize}
\end{dependencies}
